% 由 tutorial.tex 自动拆分为第 17 章
\chapter{习题与代码索引}
\begin{enumerate}[leftmargin=1.6em,itemsep=6pt]
  \item 证明：若 $q,k$ 的各分量独立且方差为 $1$，则 $\mathrm{Var}(q^\top k)=d_h$。
        若方差为 $\sigma^2$ 呢？缩放因子应改成什么？
  \item 从 $\partial a_i/\partial s_j=a_i(\delta_{ij}-a_j)$ 出发，推导
        $\partial\mathcal{L}/\partial s_j=a_j(g_j-\sum_i a_ig_i)$。
  \item 把式 \eqref{eq:lnback} 展开成 $\partial\mu,\partial\sigma^2$ 两步，验证与紧凑式等价。
  \item 去掉残差（把 $z\leftarrow z+f(z)$ 改成 $z\leftarrow f(z)$），重跑 \texttt{parity\_test.py}
        与训练，比较第 0 层梯度范数与收敛速度。
  \item 把掩码值从 $-10^9$ 改成 $-1$，看训练会发生什么，并解释原因。
  \item 用 \texttt{numpy\_from\_state()} 加载生产权重，把 $d_h$ 从 8 改成 4（保持参数量），
        重训一次，比较闭环成功率。
\end{enumerate}

\section*{代码索引}
\addcontentsline{toc}{section}{代码索引}
\begin{center}
\small
\begin{tabular}{@{}ll@{}}
\toprule
文件 & 作用 \\
\midrule
\texttt{tinytf\_numpy.py} & 纯 numpy 实现：前向 + 手写反向 + AdamW \\
\texttt{tinytf\_torch.py} & 等价的 torch 实现 + 权重互转 \\
\texttt{parity\_test.py} & 与 autograd/有限差分对齐（本教程的可信度基石）\\
\texttt{train\_demo.py} & 双实现同条件训练对比（损失曲线、耗时）\\
\texttt{load\_pretrained.py} & 加载生产权重：一致性、闭环、注意力 \\
\texttt{make\_figs.py} & 生成文中四张图 \\
\texttt{optim\_demo.py} & 优化器对比实验（toy 轨迹 + 真实模型曲线）\\
\bottomrule
\end{tabular}
\end{center}

